\documentclass[11pt,a4paper]{article}
\usepackage[utf8]{inputenc}
\usepackage[T1]{fontenc}
\usepackage{lmodern}
\usepackage[margin=1in]{geometry}
\usepackage{enumitem}
\usepackage{xcolor}
\usepackage{hyperref}
\usepackage{booktabs}
\usepackage{longtable}
\usepackage{array}
\usepackage[most]{tcolorbox}

% Define custom colors
\definecolor{promptblue}{RGB}{41, 98, 150}
\definecolor{promptbg}{RGB}{240, 247, 255}
\definecolor{promptborder}{RGB}{100, 160, 220}
\definecolor{accentgreen}{RGB}{0, 128, 100}
\definecolor{warningorange}{RGB}{200, 120, 0}

\newtcolorbox{resultbox}[1][]{
    colback=promptbg,
    colframe=promptborder,
    coltitle=white,
    fonttitle=\bfseries\small,
    title={#1},
    colbacktitle=promptblue,
    boxrule=1pt,
    arc=2mm,
    left=8pt, right=8pt, top=6pt, bottom=6pt,
    enhanced, breakable
}

\newtcolorbox{gapbox}[1][]{
    colback=orange!5,
    colframe=warningorange,
    coltitle=white,
    fonttitle=\bfseries\small,
    title={#1},
    colbacktitle=warningorange,
    boxrule=1pt,
    arc=2mm,
    left=8pt, right=8pt, top=6pt, bottom=6pt,
    enhanced, breakable
}

\hypersetup{
    colorlinks=true,
    linkcolor=promptblue,
    urlcolor=promptblue,
    citecolor=promptblue
}

\title{\textbf{Literature Review: Femomenología del color
}}
\author{Generated by Research Accelerant Agent}
\date{\today}

\begin{document}
\maketitle



\begin{resultbox}[Search Parameters]
\begin{itemize}[leftmargin=*]
    \item \textbf{Topic:} Femomenología del color

    \item \textbf{Year Range:} 2010--2026
    \item \textbf{Studies Retrieved:} 5
    \item \textbf{Citation Format:} APA
\end{itemize}
\end{resultbox}

\vspace{1em}

\textbf{\large Reviewed Studies}\par\vspace{0.5em}


\textbf{1. Teorías de la luz y el color en la época de las Luces. De Newton a Goethe}

\textbf{Full Citation:} J Pimentel (2015). Teorías de la luz y el color en la época de las Luces. De Newton a Goethe. 	extit{Arbor}. https://doi.org/N/A

\begin{resultbox}[Study Details]
\begin{itemize}[leftmargin=*]
    \item \textbf{Authors:} J Pimentel
    \item \textbf{Year:} 2015
    \item \textbf{Journal:} Arbor
    
    
    
    \item \textbf{DOI:} \href{https://doi.org/}{N/A}
    \item \textbf{Citation Count:} 29 (Google Scholar)
    \item \textbf{Access:} \href{https://arbor.revistas.csic.es/index.php/arbor/article/view/2067}{Semantic Scholar}
\end{itemize}
\end{resultbox}

\textbf{Abstract:} … por revisar la teoría de la luz de Isaac Newton, centrándonos … entendida la luz, así como el papel simbólico de la luz, metáfora … entender mejor no sólo qué son la luz, el ojo o los fenóme…\par\vspace{0.5em}



\textbf{2. As cores ea Visão ea Visão das Cores}

\textbf{Full Citation:} MCS Ribeiro (2011). As cores ea Visão ea Visão das Cores. 	extit{2011}. https://doi.org/N/A

\begin{resultbox}[Study Details]
\begin{itemize}[leftmargin=*]
    \item \textbf{Authors:} MCS Ribeiro
    \item \textbf{Year:} 2011
    \item \textbf{Journal:} 2011
    
    
    
    \item \textbf{DOI:} \href{https://doi.org/}{N/A}
    \item \textbf{Citation Count:} 10 (Google Scholar)
    \item \textbf{Access:} \href{https://search.proquest.com/openview/d00872f699d42bae8e5936ff65e42d0a/1?pq-origsite=gscholar&cbl=2026366&diss=y}{Semantic Scholar}
\end{itemize}
\end{resultbox}

\textbf{Abstract:} … A luz é considerada como a energia que sustenta a vida. Toda … Sabemos que através da luz recebemos grandes quantidades de … aos diferentes comprimentos de onda da luz recebida. …\par\vspace{0.5em}



\textbf{3. Evaluación de las estrategias de adaptación a disfunciones de la visión del color}

\textbf{Full Citation:} I Coca Torrents (2012). Evaluación de las estrategias de adaptación a disfunciones de la visión del color. 	extit{2012}. https://doi.org/N/A

\begin{resultbox}[Study Details]
\begin{itemize}[leftmargin=*]
    \item \textbf{Authors:} I Coca Torrents
    \item \textbf{Year:} 2012
    \item \textbf{Journal:} 2012
    
    
    
    \item \textbf{DOI:} \href{https://doi.org/}{N/A}
    \item \textbf{Citation Count:} 3 (Google Scholar)
    \item \textbf{Access:} \href{https://upcommons.upc.edu/entities/publication/286e34af-fb41-42b1-8d35-5fc558f4e080}{Semantic Scholar}
\end{itemize}
\end{resultbox}

\textbf{Abstract:} … nos ayuda, por ejemplo, a saber si una fruta está lo suficientemente madura para poder comerla, aun sabiendo que el color de dicha fruta puede cambiar dependiendo de la luz que le …\par\vspace{0.5em}



\textbf{4. La visión del color: mecanismos fisiológicos, teorías y deficiencias comunes}

\textbf{Full Citation:} MS Larraz, AC Martínez, BM Lacámara… (2025). La visión del color: mecanismos fisiológicos, teorías y deficiencias comunes. 	extit{Revista Sanitaria de …}. https://doi.org/N/A

\begin{resultbox}[Study Details]
\begin{itemize}[leftmargin=*]
    \item \textbf{Authors:} MS Larraz, AC Martínez, BM Lacámara…
    \item \textbf{Year:} 2025
    \item \textbf{Journal:} Revista Sanitaria de …
    
    
    
    \item \textbf{DOI:} \href{https://doi.org/}{N/A}
    \item \textbf{Citation Count:} 0 (Google Scholar)
    \item \textbf{Access:} \href{https://dialnet.unirioja.es/servlet/articulo?codigo=10349656}{Semantic Scholar}
\end{itemize}
\end{resultbox}

\textbf{Abstract:} … en color) y bastones (para la visión nocturna). Los conos están distribuidos en la fóvea y responden a las longitudes de onda de luz … In conclusion, color vision is fundamental both in …\par\vspace{0.5em}



\textbf{5. LA PERCEPCIÓN DE LA LUZ POR EL SISTEMA VISUAL HUMANO. EL COLOR}

\textbf{Full Citation:} JR Mora (2026). LA PERCEPCIÓN DE LA LUZ POR EL SISTEMA VISUAL HUMANO. EL COLOR. 	extit{Luz y Vida.}. https://doi.org/N/A

\begin{resultbox}[Study Details]
\begin{itemize}[leftmargin=*]
    \item \textbf{Authors:} JR Mora
    \item \textbf{Year:} 2026
    \item \textbf{Journal:} Luz y Vida.
    
    
    
    \item \textbf{DOI:} \href{https://doi.org/}{N/A}
    \item \textbf{Citation Count:} 0 (Google Scholar)
    \item \textbf{Access:} \href{https://www.rasc.es/wp-content/uploads/2025/01/Luz-y-Vida_Conmemorando-el-dia-internacional-de-la-luz.pdf#page=89}{Semantic Scholar}
\end{itemize}
\end{resultbox}

\textbf{Abstract:} … absoluta de luz y la saturación relacionada con la pureza del color que presente la luz en … sistema de representación del color que ayudaran a medir el color con números. Esto era …\par\vspace{0.5em}



\vspace{1em}
\textbf{\large Summary Statistics}\par\vspace{0.5em}

\begin{center}
\begin{tabular}{lc}
\toprule
\textbf{Metric} & \textbf{Value} \\
\midrule
Total Studies Reviewed & 5 \\
Average Citations per Study & 8 \\
Publication Year Range & 2011--2026 \\
Studies with Open Access PDF & 4 \\
\bottomrule
\end{tabular}
\end{center}


\vfill
\begin{center}
\textit{Document generated by Research Accelerant Agent on \today.}
\end{center}

\end{document}